\newpage

% ######################################################################
% PART XVII: BACKEND FILES + REVISION PLAN
% ######################################################################
\part{Backend Files + Daily Revision Plan}

% ######################################################################
% SECTION 67: REMAINING BACKEND FILES
% ######################################################################
\section{Remaining Backend Files --- Serializers, URLs, Conftest}

\subsection{serializers.py --- Full Walkthrough}

\begin{codestorybox}
\begin{lstlisting}[style=pythonstyle, caption={serializers.py --- dono serializers}]
from rest_framework import serializers
from .models import ResearchSession, BenchmarkEvaluation

class ResearchSessionSerializer(serializers.ModelSerializer):
    class Meta:
        model = ResearchSession
        fields = ["id", "topic", "report", "messages",
                  "sources_count", "revision_count", "created_at"]

class BenchmarkEvaluationSerializer(serializers.ModelSerializer):
    class Meta:
        model = BenchmarkEvaluation
        fields = ["id", "topic", "single_agent_report",
                  "multi_agent_report", "single_agent_depth",
                  "multi_agent_depth", "single_agent_verifiability",
                  "multi_agent_verifiability", "verdict", "created_at"]
\end{lstlisting}

\textbf{Design:} (1) \texttt{ModelSerializer} --- model fields se auto field mapping + validation; (2) Explicit \texttt{fields} list --- kya expose karna hai control (user field exclude); (3) \texttt{user} intentionally serializer mein nahi --- request se attach hota hai, client se nahi. Interview point: ``Serializer = API ka contract --- kya dikhega, kya validation, kya chhupa hoga, sab yahan control hota hai.''
\end{codestorybox}

\subsection{urls.py --- Full Walkthrough}

\begin{codestorybox}
\begin{lstlisting}[style=pythonstyle, caption={api/urls.py}]
from django.urls import path
from . import views

urlpatterns = [
    path("health/", views.HealthCheckView.as_view()),
    path("auth/register/", views.RegisterView.as_view()),
    path("auth/login/", views.LoginView.as_view()),
    path("research/execute/", views.ExecuteResearchAPIView.as_view()),
    path("research/benchmark/", views.BenchmarkAPIView.as_view()),
    path("benchmark/history/", views.BenchmarkHistoryAPIView.as_view()),
]
\end{lstlisting}

\textbf{Routing structure:} (1) \texttt{health/} top-level; (2) \texttt{auth/} namespace --- register/login; (3) \texttt{research/} namespace --- execute/benchmark; (4) \texttt{benchmark/history/} read-only. Ye REST-consistent nested design hai. Main urls.py mein \texttt{path("api/", include("api.urls"))} --- prefix ek jagah.
\end{codestorybox}

\subsection{conftest.py --- Mocking Infrastructure}

\begin{codestorybox}
\begin{lstlisting}[style=pythonstyle, caption={conftest.py --- fixtures}]
import pytest
from unittest.mock import patch, MagicMock

@pytest.fixture
def mock_research_graph():
    mock = MagicMock()
    mock.invoke.return_value = {
        "final_report": "# Report\n\nTest content",
        "messages": ["Agent ran"],
    }
    with patch("graph.workflow.research_graph", mock):
        yield mock

@pytest.fixture
def exploding_research_graph():
    mock = MagicMock()
    mock.invoke.side_effect = RuntimeError("pipeline crashed")
    with patch("graph.workflow.research_graph", mock):
        yield mock

@pytest.fixture
def mock_benchmark():
    with patch("benchmark.evaluator.run_single_agent_baseline",
               return_value="# Baseline report"),
         patch("benchmark.evaluator.evaluate_outputs",
               return_value={"single_agent_depth": 9, ...}):
        yield
\end{lstlisting}

\textbf{Fixture design:} (1) \texttt{MagicMock} --- auto attributes/methods; (2) \texttt{side\_effect} --- exception simulation (error paths); (3) \texttt{return\_value} --- success paths; (4) \texttt{patch} context manager --- test ke baad automatic restore. Interview point: ``Fixtures = test setup ka reusable building block. Mocking se tests fast (9s), free, aur deterministic hote hain --- real API calls sirf manual E2E mein.''
\end{codestorybox}

\subsection{Settings/wsgi/asgi/manage.py --- Quick Tour}

\begin{codestorybox}
\begin{lstlisting}[style=bashstyle, caption={Django entry files}]
# manage.py: CLI entry (runserver, migrate, shell)
python manage.py runserver 8001

# settings.py: INSTALLED_APPS (rest_framework, corsheaders, api),
# DATABASES (sqlite3), CORS_ALLOWED_ORIGINS, rest_framework auth classes

# wsgi.py: sync WSGI app (Gunicorn production ke liye)
# asgi.py: async ASGI app (uvicorn / real-time ke liye)
\end{lstlisting}

\textbf{Interview points:} (1) WSGI vs ASGI --- sync vs async server interface; (2) INSTALLED\_APPS = Django ka plugin system; (3) CORS config --- Streamlit origin allow; (4) REST framework settings --- DefaultAuthentication (token/session).
\end{codestorybox}

% ######################################################################
% SECTION 68: DAILY REVISION PLAN
% ######################################################################
\section{Daily Revision Plan --- Interview Tak}

\begin{longtable}{p{3.2cm} p{10.8cm}}
\toprule
\textbf{Day} & \textbf{Revision focus} \\
\midrule
Day 1 & Parts I-III --- setup, concepts, code walkthrough (2x speed read) \\
Day 2 & Part IV --- deep dives + prompts verbatim + system design \\
Day 3 & Part V --- 68 Q\&A (bol ke practice) + pitches + whiteboard \\
Day 4 & Parts VI-VII --- API reference, tests, deployment, glossary \\
Day 5 & Parts VIII-IX --- transformer theory, debugging stories, 30 more Q\&A \\
Day 6 & Parts X-XI --- DSA patterns, core CS, demo scripts, mock interview \\
Day 7 & Parts XII-XVI --- Python deep dive, quick-fire, cheat card, full mock \\
Daily & 30-min tutorial script 1x + 10 numbers + 4 stories \\
\bottomrule
\end{longtable}

\begin{memorybox}
\textbf{Interview day checklist:}
\begin{enumerate}[label=\arabic*.]
\item Numbers yaad karo (cheat card section 66).
\item 30-sec pitch mirror ke saamne 2 baar.
\item Ek mock interview kisi dost ke saath.
\item Laptop ready: code + demo + tests chalne ready.
\item Setup: local server ports free, API keys loaded, internet on.
\item Deep breath. \textbf{Tum taiyaar ho.}
\end{enumerate}
\end{memorybox}

\begin{storybox}
\textbf{17 parts, 68 sections, 150+ pages.} Har logic explain kiya, har code line samjhaya, har interview sawaal ka jawab diya --- sab Hinglish mein, story style mein. Is document ne aapko \textbf{gyaan} diya hai; ab \textbf{confident} bolo. Interview sirf ek conversation hai jisme aap apna kaam dikha rahe ho. \textbf{All the best!} \faGraduationCap
\end{storybox}
